\documentclass[12pt,a4paper]{article}

% ── Packages ──────────────────────────────────────────────────────────────────
\usepackage[T1]{fontenc}
\usepackage[utf8]{inputenc}
\usepackage{lmodern}
\usepackage[margin=1in]{geometry}
\usepackage{graphicx}
\usepackage{amsmath}
\usepackage{booktabs}
\usepackage{array}
\usepackage{tabularx}
\usepackage{longtable}
\usepackage{listings}
\usepackage{xcolor}
\usepackage{hyperref}
\usepackage{fancyhdr}
\usepackage{titlesec}
\usepackage{parskip}
\usepackage{enumitem}
\usepackage{microtype}
\usepackage{abstract}
\usepackage{mdframed}
\usepackage{float}

% ── Colours ───────────────────────────────────────────────────────────────────
\definecolor{navyblue}{HTML}{1e2a45}
\definecolor{accentblue}{HTML}{2563eb}
\definecolor{lightgray}{HTML}{f1f5f9}
\definecolor{midgray}{HTML}{64748b}
\definecolor{darkgray}{HTML}{1e293b}
\definecolor{codeblue}{HTML}{1d4ed8}
\definecolor{codebg}{HTML}{f8fafc}
\definecolor{borderblue}{HTML}{bfdbfe}
\definecolor{successgreen}{HTML}{16a34a}
\definecolor{warnred}{HTML}{dc2626}

% ── Hyperlinks ────────────────────────────────────────────────────────────────
\hypersetup{
  colorlinks=true,
  linkcolor=accentblue,
  urlcolor=accentblue,
  citecolor=accentblue,
  pdftitle={Warehouse Inventory Management System -- Project Report},
  pdfauthor={Bhavith Chandra},
  pdfsubject={CS6103 Final Project},
}

% ── Section styling ───────────────────────────────────────────────────────────
\titleformat{\section}
  {\large\bfseries\color{navyblue}}
  {\thesection.}{0.6em}{}
  [\color{accentblue}\titlerule]

\titleformat{\subsection}
  {\normalsize\bfseries\color{darkgray}}
  {\thesubsection}{0.6em}{}

\titlespacing{\section}{0pt}{18pt}{8pt}
\titlespacing{\subsection}{0pt}{10pt}{4pt}

% ── Listings (code blocks) ────────────────────────────────────────────────────
\lstdefinestyle{javastyle}{
  language=Java,
  basicstyle=\ttfamily\small,
  keywordstyle=\color{accentblue}\bfseries,
  commentstyle=\color{midgray}\itshape,
  stringstyle=\color{successgreen},
  numberstyle=\tiny\color{midgray},
  backgroundcolor=\color{codebg},
  frame=single,
  framerule=0.4pt,
  rulecolor=\color{borderblue},
  numbers=left,
  numbersep=8pt,
  breaklines=true,
  captionpos=b,
  showstringspaces=false,
  tabsize=4,
  xleftmargin=14pt,
}

\lstdefinestyle{shell}{
  basicstyle=\ttfamily\small,
  backgroundcolor=\color{codebg},
  frame=single,
  framerule=0.4pt,
  rulecolor=\color{borderblue},
  breaklines=true,
  captionpos=b,
  xleftmargin=14pt,
}

\lstset{style=javastyle}

% ── Header / Footer ───────────────────────────────────────────────────────────
\pagestyle{fancy}
\fancyhf{}
\renewcommand{\headrulewidth}{0.4pt}
\renewcommand{\headrule}{\hbox to\headwidth{\color{accentblue}\leaders\hrule height \headrulewidth\hfill}}
\fancyhead[L]{\small\color{midgray}Warehouse Inventory Management System}
\fancyhead[R]{\small\color{midgray}CS6103 --- Final Project Report}
\fancyfoot[C]{\small\color{midgray}\thepage}

% ── Info box ──────────────────────────────────────────────────────────────────
\newmdenv[
  backgroundcolor=codebg,
  linecolor=accentblue,
  linewidth=1pt,
  leftmargin=0pt,
  rightmargin=0pt,
  innerleftmargin=10pt,
  innerrightmargin=10pt,
  innertopmargin=8pt,
  innerbottommargin=8pt,
]{infobox}

% ── Result box ────────────────────────────────────────────────────────────────
\newmdenv[
  backgroundcolor=codebg,
  linecolor=successgreen,
  linewidth=1.5pt,
  leftline=true,
  rightline=false,
  topline=false,
  bottomline=false,
  leftmargin=0pt,
  rightmargin=0pt,
  innerleftmargin=12pt,
  innerrightmargin=10pt,
  innertopmargin=6pt,
  innerbottommargin=6pt,
]{resultbox}

% ── Document ──────────────────────────────────────────────────────────────────
\begin{document}

% ── Title page ────────────────────────────────────────────────────────────────
\begin{titlepage}
  \centering
  \vspace*{2cm}

  {\Huge\bfseries\color{navyblue} Warehouse Inventory\\[6pt] Management System\par}
  \vspace{0.4cm}
  {\large\color{accentblue}\rule{0.6\linewidth}{1.5pt}\par}
  \vspace{0.4cm}
  {\large\color{midgray} Final Project Report\par}

  \vspace{2.5cm}

  \begin{center}
    \begin{tabular}{rl}
      \textbf{Course}    & CS6103 --- Introduction to Java \\[4pt]
      \textbf{Semester}  & Spring 2026 \\[4pt]
      \textbf{Student}   & Bhavith Chandra \\[4pt]
      \textbf{Net ID}    & bc4066 \\[4pt]
      \textbf{Repository}& \href{https://github.com/Bhavith-Chandra/warehouse-inventory-management}%
                              {\texttt{github.com/Bhavith-Chandra/warehouse-inventory-management}}
    \end{tabular}
  \end{center}

  \vfill
  {\small\color{midgray}\today}
\end{titlepage}

% ── Table of Contents ─────────────────────────────────────────────────────────
\tableofcontents
\newpage

% ── Abstract ─────────────────────────────────────────────────────────────────
\begin{abstract}
This report presents the design and implementation of a multi-user Warehouse Inventory
Management System developed as a final project for CS6103. The system follows a
three-tier client-server architecture built entirely in Java, using TCP sockets for
networking, JavaFX~21 for the graphical user interface, and SQLite for persistent
storage. A central focus of the implementation is correctness under concurrency: the
system guarantees that two clients attempting to deduct the last unit of a product
simultaneously will never produce a double-deduction---a property proven by an automated
stress test included with the submission. All 29 automated end-to-end verification checks
pass.
\end{abstract}

% ─────────────────────────────────────────────────────────────────────────────
\section{Introduction}
% ─────────────────────────────────────────────────────────────────────────────

Modern warehouses require real-time, multi-user access to inventory data. Staff need to
log stock movements instantly, managers need accurate stock levels at a glance, and the
system must remain consistent even when multiple users act simultaneously. This project
builds a fully functional desktop application that satisfies all of these requirements
using the core Java technologies covered in CS6103.

The system delivers the following capabilities:

\begin{itemize}[leftmargin=1.4em]
  \item Real-time product catalogue management (create, read, update, delete)
  \item Stock-level tracking via add, remove, and absolute-adjust operations
  \item Automatic low-stock alerting with inline restock shortcuts
  \item A full, immutable audit transaction log with timestamps and performer tracking
  \item Role-based access control (\texttt{ADMIN} vs.\ \texttt{STAFF})
  \item Live push-event updates so all connected clients refresh automatically
  \item CSV export of both inventory and transaction history
\end{itemize}

% ─────────────────────────────────────────────────────────────────────────────
\section{System Architecture}
% ─────────────────────────────────────────────────────────────────────────────

The application is structured as a strict three-tier architecture that isolates
presentation, business logic, and persistence.

\subsection{Tier 1 --- Presentation (JavaFX Client)}

The client is a JavaFX~21 desktop application (\texttt{InventoryClient.java}).
It contains no business logic; every user action is serialised into a
\texttt{Request} object and transmitted over a TCP socket. The UI runs entirely
on the JavaFX Application Thread; all network I/O runs on daemon background
threads, and results are marshalled back to the UI via
\texttt{Platform.runLater()}.

\subsection{Tier 2 --- Application Server (Java TCP Server)}

\texttt{InventoryServer} opens a \texttt{ServerSocket} on port~5555 and spawns
one \texttt{Thread(ClientHandler)} per accepted connection. \texttt{ClientHandler}
reads \texttt{Request} objects, dispatches them to \texttt{InventoryService} (the
business-logic layer), and writes \texttt{Response} objects back. After every
successful mutation the server broadcasts a push event to all subscribed clients so
their UIs refresh without polling.

\subsection{Tier 3 --- Data Layer (SQLite via JDBC)}

\texttt{DatabaseManager} owns a single JDBC \texttt{Connection} to an SQLite file
(\texttt{data/warehouse.db}). WAL (Write-Ahead Logging) journal mode is enabled so
concurrent reads never block an in-progress write. The schema comprises three tables:
\texttt{users}, \texttt{products}, and \texttt{transactions}.

\subsection{Wire Protocol}

All messages are Java-serialised objects. \texttt{Request} carries a \texttt{Type}
enum (14 verbs: \texttt{LOGIN}, \texttt{LIST\_PRODUCTS}, \texttt{ADD\_STOCK},
\texttt{REMOVE\_STOCK}, etc.), a parameter \texttt{Map<String,Object>}, and a
\texttt{requestId = System.nanoTime()}. \texttt{Response} carries a \texttt{Status}
enum (\texttt{OK} / \texttt{ERROR} / \texttt{EVENT}), a payload, and the matching
\texttt{requestId}. The client's \texttt{ServerConnection} correlates incoming
responses to outstanding \texttt{CompletableFuture}s by \texttt{requestId},
enabling fully asynchronous networking over a single socket without blocking the UI
thread.

\begin{figure}[H]
\centering
\begin{infobox}
\begin{lstlisting}[style=shell,numbers=none,basicstyle=\ttfamily\footnotesize]
+------------------+    TCP / Serialisation    +------------------+
|  JavaFX Client   | <-----------------------> | InventoryServer  |
| InventoryClient  |                           |   (port 5555)    |
+------------------+                           +--------+---------+
                                                        | one Thread per client
                                          +-------------+-------------+
                                     ClientHandler  ClientHandler  ClientHandler
                                          |
                                 +--------v---------+
                                 | InventoryService | ReentrantLock + JDBC tx
                                 +--------+---------+
                                          |
                                 +--------v---------+
                                 |  DatabaseManager | SQLite WAL
                                 +------------------+
\end{lstlisting}
\end{infobox}
\caption{High-level architecture diagram}
\end{figure}

% ─────────────────────────────────────────────────────────────────────────────
\section{Concurrency Design}
% ─────────────────────────────────────────────────────────────────────────────

The headline correctness property from the project proposal is:

\begin{quote}
\textit{``If two users try to deduct the last unit of a product at the same time,
only one should succeed.''}
\end{quote}

\texttt{InventoryService.adjustStock()} enforces this with a two-layer defence.

\subsection{Layer 1 --- Per-Product \texttt{ReentrantLock}}

A \texttt{ConcurrentHashMap<Integer, ReentrantLock>} maps each product ID to its
own lock. Every stock mutation acquires the lock for \emph{that} product before
performing any I/O. Writes to different products proceed in parallel; writes to the
\emph{same} product are strictly serialised within the JVM.

\subsection{Layer 2 --- JDBC Transaction}

Inside the lock the code:

\begin{enumerate}[leftmargin=1.6em]
  \item Calls \texttt{setAutoCommit(false)} to open a transaction.
  \item Re-reads the current quantity from the database \emph{inside the same
        transaction} (preventing stale cached values from leaking through).
  \item Validates the business rule: \texttt{quantity + delta >= 0}; throws
        \texttt{InsufficientStockException} otherwise.
  \item Applies the \texttt{UPDATE products SET quantity = ?}.
  \item Writes an audit row to \texttt{transactions}.
  \item Calls \texttt{commit()}. On any exception, \texttt{safeRollback()} is called.
\end{enumerate}

The database schema also enforces \texttt{CHECK (quantity >= 0)} as a final
safety net at the storage layer.

\subsection{Why Two Layers?}

The database constraint alone would allow both clients to \emph{attempt} the
deduction; one would receive a confusing \texttt{SQLException} from the
\texttt{CHECK} violation. The application-level \texttt{ReentrantLock} ensures
that exactly one client gets a clean \textit{``Insufficient stock''} message and
the other a successful response---the behaviour users and downstream systems
expect.

\begin{lstlisting}[style=javastyle, caption={Core of \texttt{adjustStock()} --- two-layer concurrency guard}]
ReentrantLock lock = lockFor(productId);
lock.lock();
try {
    synchronized (db) {
        Connection c = db.getConnection();
        c.setAutoCommit(false);
        try {
            // Re-read inside transaction to prevent stale reads
            int current = readQuantity(c, productId);
            int next    = current + delta;
            if (next < 0)
                throw new InsufficientStockException(...);

            updateQuantity(c, productId, next);
            writeAuditLog(c, productId, delta, next, performedBy, reason);
            c.commit();
        } catch (Exception e) {
            safeRollback(c);
            throw e;
        }
    }
} finally {
    lock.unlock();
}
\end{lstlisting}

% ─────────────────────────────────────────────────────────────────────────────
\section{Authentication and Security}
% ─────────────────────────────────────────────────────────────────────────────

\subsection{Password Storage}

Passwords are stored as salted SHA-256 hashes using \texttt{PasswordHasher.java}.
The on-disk format is:

\begin{center}
  \texttt{base64(salt) : base64(SHA-256(salt || password))}
\end{center}

Constant-time comparison (\texttt{MessageDigest.isEqual}) is used during
verification to prevent timing-oracle attacks.

\subsection{Role-Based Access Control}

\begin{table}[H]
\centering
\begin{tabular}{lll}
\toprule
\textbf{Role} & \textbf{Permitted Operations} & \textbf{Denied Operations} \\
\midrule
\texttt{ADMIN} & Full access (CRUD + stock + log) & --- \\
\texttt{STAFF} & Stock operations, read-only views & Create/update/delete products \\
\bottomrule
\end{tabular}
\caption{Role permissions}
\end{table}

Role checks are enforced in \texttt{ClientHandler.requireRole()} before every
admin-only operation. The server also enforces a login-first protocol: any request
other than \texttt{LOGIN} before authentication is rejected with
\textit{``Not authenticated; LOGIN required first.''}.

\subsection{Default Credentials (First Run)}

\begin{table}[H]
\centering
\begin{tabular}{lll}
\toprule
\textbf{Username} & \textbf{Password} & \textbf{Role} \\
\midrule
\texttt{admin} & \texttt{admin123} & \texttt{ADMIN} \\
\texttt{staff} & \texttt{staff123} & \texttt{STAFF} \\
\bottomrule
\end{tabular}
\caption{Seeded default credentials}
\end{table}

% ─────────────────────────────────────────────────────────────────────────────
\section{Implementation Overview}
% ─────────────────────────────────────────────────────────────────────────────

\subsection{Server-Side Classes}

\begin{table}[H]
\centering
\begin{tabularx}{\linewidth}{lX}
\toprule
\textbf{Class} & \textbf{Responsibility} \\
\midrule
\texttt{InventoryServer}    & \texttt{ServerSocket} accept loop; handler registry; broadcast engine \\
\texttt{ClientHandler}      & Per-connection request dispatch; \texttt{writeLock} for thread-safe socket writes \\
\texttt{InventoryService}   & All business logic; concurrency guards; JDBC transaction management \\
\texttt{DatabaseManager}    & Schema creation; WAL/FK PRAGMAs; single JDBC connection lifetime \\
\bottomrule
\end{tabularx}
\caption{Server-side classes}
\end{table}

\subsection{Client-Side Classes}

\begin{table}[H]
\centering
\begin{tabularx}{\linewidth}{lX}
\toprule
\textbf{Class} & \textbf{Responsibility} \\
\midrule
\texttt{InventoryClient}    & JavaFX Application; sidebar navigation; keyboard shortcuts \\
\texttt{ServerConnection}   & Async socket I/O; CompletableFuture correlation; event dispatch \\
\texttt{DashboardView}      & KPI cards, PieChart by category, BarChart top-10, activity feed \\
\texttt{InventoryView}      & Searchable/filterable TableView; Add/Remove Stock dialogs \\
\texttt{StockOpsView}       & Dedicated receive/ship/cycle-count panel with live stock preview \\
\texttt{LowStockView}       & Auto-filtered low-stock list with one-click Restock button \\
\texttt{HistoryView}        & Transaction audit log; colour-coded action badges; CSV export \\
\bottomrule
\end{tabularx}
\caption{Client-side classes}
\end{table}

\subsection{Shared / Utility Classes}

\begin{itemize}[leftmargin=1.4em]
  \item \texttt{model/} --- \texttt{Product}, \texttt{TransactionLog}, \texttt{User}
        (\texttt{Serializable} domain objects transported over the wire)
  \item \texttt{protocol/} --- \texttt{Request} (14-verb \texttt{Type} enum),
        \texttt{Response} (\texttt{Status} + \texttt{EventType} enums)
  \item \texttt{util/PasswordHasher} --- salted SHA-256 hash and verify
\end{itemize}

% ─────────────────────────────────────────────────────────────────────────────
\section{Testing and Verification}
% ─────────────────────────────────────────────────────────────────────────────

Two automated test programs are included under \texttt{src/com/warehouse/test/}.

\subsection{\texttt{ProposalVerify.java} --- End-to-End Verification}

Connects to a live server and exercises every project-proposal requirement through
29 automated checks covering:

\begin{itemize}[leftmargin=1.4em]
  \item All six CRUD operations (\texttt{ADD\_PRODUCT}, \texttt{LIST\_PRODUCTS},
        \texttt{SEARCH\_PRODUCTS}, \texttt{GET\_PRODUCT}, \texttt{UPDATE\_PRODUCT},
        \texttt{REMOVE\_PRODUCT})
  \item All three stock-mutation types (add, remove, absolute adjust)
  \item Over-deduction guard (insufficient-stock rejection)
  \item Low-stock report inclusion
  \item Full transaction-log inspection (type, performer, timestamp)
  \item Role-based access-control enforcement
  \item Unauthenticated-request rejection
  \item Push-event delivery to subscribed clients
\end{itemize}

\begin{resultbox}
\textbf{Result:}\quad \textcolor{successgreen}{\textbf{29 / 29 PASS}}
\end{resultbox}

\subsection{\texttt{ConcurrencyStressTest.java} --- Race-Condition Proof}

Resets a chosen product to $N$ units, launches $T$ client threads each
repeatedly calling \texttt{REMOVE\_STOCK(1)} until rejected, then reads the
final stock level.

\begin{lstlisting}[style=shell, numbers=none,
  caption={Sample stress-test output ($N=50$, $T=15$)}]
=== Warehouse Concurrency Stress Test ===
productId=1  startStock=50  threads=15

  successful deductions:   50
  'insufficient' rejects:  15
  other errors:             0
  final stock:              0

PASS: exactly 50 units deducted, no double-deduction, final stock 0.
\end{lstlisting}

The test confirms that under heavy parallel contention the invariant
$\textit{successes} = N \;\wedge\; \textit{finalStock} = 0$
always holds, and that no double-deduction ever occurs.

% ─────────────────────────────────────────────────────────────────────────────
\section{How to Run}
% ─────────────────────────────────────────────────────────────────────────────

\subsection{Prerequisites}

\begin{itemize}[leftmargin=1.4em]
  \item JDK~17 or later (tested on OpenJDK~25)
  \item JavaFX SDK~21.0.5 --- download from
        \url{https://gluonhq.com/products/javafx/} and extract to
        \texttt{lib/javafx-sdk-21.0.5/}
  \item \texttt{sqlite-jdbc-3.46.1.3.jar}, \texttt{slf4j-api-2.0.13.jar},
        and \texttt{slf4j-simple-2.0.13.jar} are already included in
        \texttt{lib/}
\end{itemize}

\subsection{Build and Launch}

\begin{lstlisting}[style=shell, numbers=none, caption={Step 1 --- Compile}]
bash scripts/compile.sh
\end{lstlisting}

\begin{lstlisting}[style=shell, numbers=none, caption={Step 2 --- Start the server (Terminal 1)}]
bash scripts/run-server.sh
# First run: creates data/warehouse.db, seeds 2 users + 10 sample products
# Output: "Warehouse server listening on port 5555"
\end{lstlisting}

\begin{lstlisting}[style=shell, numbers=none, caption={Step 3 --- Start the client (Terminal 2, repeat for multi-user demo)}]
bash scripts/run-client.sh
# Connection dialog: localhost:5555 / admin / admin123
\end{lstlisting}

\begin{lstlisting}[style=shell, numbers=none, caption={Step 4 (optional) --- Concurrency stress test}]
bash scripts/run-stress.sh
# Custom: bash scripts/run-stress.sh [productId] [startStock] [threads]
\end{lstlisting}

\subsection{Eclipse IDE Import}

\begin{enumerate}[leftmargin=1.6em]
  \item \textbf{File} $\to$ \textbf{Import} $\to$ \textbf{General} $\to$
        \textbf{Existing Projects into Workspace} $\to$ select this folder.
  \item The included \texttt{.project} and \texttt{.classpath} configure all
        JAR dependencies automatically.
  \item Add the following VM arguments to the client run configuration:
\end{enumerate}

\begin{lstlisting}[style=shell, numbers=none, basicstyle=\ttfamily\small]
--module-path lib/javafx-sdk-21.0.5/lib
--add-modules javafx.controls,javafx.fxml,javafx.graphics,javafx.base
\end{lstlisting}

\subsection{UI Navigation and Keyboard Shortcuts}

\begin{table}[H]
\centering
\begin{tabular}{lll}
\toprule
\textbf{View / Action} & \textbf{Description} & \textbf{Shortcut} \\
\midrule
Dashboard           & KPI cards, charts, live activity feed               & --- \\
Inventory           & Product table; search, filter, CRUD, stock dialogs  & Ctrl+F / Ctrl+N \\
Stock Operations    & Receive / ship / cycle-count panel                  & --- \\
Low Stock Alerts    & Products at or below reorder level                  & --- \\
Transaction History & Full audit log with colour-coded badges             & --- \\
Refresh all         & Pull latest data from server                        & Ctrl+R \\
Export CSV          & Save current inventory snapshot to CSV              & Ctrl+E \\
\bottomrule
\end{tabular}
\caption{UI navigation and shortcuts}
\end{table}

% ─────────────────────────────────────────────────────────────────────────────
\section{Mapping to CS6103 Advanced Java Concepts}
% ─────────────────────────────────────────────────────────────────────────────

\begin{table}[H]
\centering
\begin{tabularx}{\linewidth}{lX}
\toprule
\textbf{Concept} & \textbf{Where It Appears} \\
\midrule
JavaFX GUI              & \texttt{InventoryClient} --- TableView, PieChart, BarChart, Dialog, FilteredList, CSS theming \\
Java Sockets            & \texttt{InventoryServer} (ServerSocket), \texttt{ServerConnection} (Socket, ObjectIn/OutputStream) \\
Multithreading          & Accept loop spawns one Thread(ClientHandler) per client; Platform.runLater() for UI marshalling \\
JDBC + concurrency      & \texttt{DatabaseManager} (WAL, schema); \texttt{InventoryService} (ReentrantLock, setAutoCommit, commit/rollback) \\
\midrule
\multicolumn{2}{l}{\textit{Beyond the proposal}} \\
\midrule
Push-event broadcast    & Server broadcasts \texttt{Response.event()} to all subscribed clients after every mutation \\
Salted password hashing & \texttt{PasswordHasher} --- SHA-256 with per-user random salt, constant-time verify \\
Automated testing       & \texttt{ProposalVerify} (29 checks), \texttt{ConcurrencyStressTest} (race proof) \\
\bottomrule
\end{tabularx}
\caption{Project proposal concepts mapped to implementation}
\end{table}

% ─────────────────────────────────────────────────────────────────────────────
\section{Conclusion}
% ─────────────────────────────────────────────────────────────────────────────

The Warehouse Inventory Management System demonstrates the application of core Java
concepts---TCP sockets, multithreading, JavaFX, and JDBC---in a cohesive,
production-quality application. The two-layer concurrency strategy (per-product
\texttt{ReentrantLock} combined with a JDBC transaction and a re-read of the
current stock inside that transaction) correctly solves the classic ``last unit''
race condition identified in the project proposal as the central correctness
challenge.

All 29 automated proposal-verification checks pass, and the concurrency stress test
confirms that no double-deduction occurs under heavy parallel load across 15
concurrent client threads.

The full source code, scripts, and this report are available at:\\[4pt]
\url{https://github.com/Bhavith-Chandra/warehouse-inventory-management}

\end{document}
